\documentclass[11pt,a4paper]{article}

\usepackage[margin=1in]{geometry}
\usepackage{amsmath,amssymb,amsthm}
\usepackage{microtype}
\usepackage[round]{natbib}
\usepackage{xcolor}
\usepackage{hyperref}
\hypersetup{colorlinks=true,linkcolor=blue!50!black,citecolor=blue!50!black,urlcolor=blue!50!black}
\usepackage{enumitem}
\setlist{topsep=2pt,itemsep=2pt,parsep=0pt}

\title{Demand-Based Asset Pricing:\\
\large A Note on the Framework, Non-Homotheticity, and Negative (Net) Portfolio Shares}
\author{}
\date{\today}

\newcommand{\eps}{\varepsilon}
\newcommand{\E}{\mathbb{E}}
\newcommand{\me}{\mathit{me}}

\begin{document}
\maketitle

\section{Introduction}

Demand-based asset pricing reframes the canonical asset-pricing question.
Rather than deriving prices from a representative-agent stochastic
discount factor and asking which preferences and shocks rationalise the
cross-section, it specifies an estimable system of investor-level demand
functions over assets, identifies its parameters from observed
holdings, and lets prices clear the resulting demand against fixed
supply. The intellectual lineage runs from McFadden's discrete-choice
demand \citep{mcfadden1974conditional}, to Berry's inversion approach
\citep{berry1994estimating} and the BLP random-coefficients framework
\citep{blp1995,nevo2001measuring}, to \citet{koijen2019demand} who first
ported these tools to portfolio choice. The framework has since
underpinned work on inelastic markets
\citep{gabaix2021inelastic}, central-bank asset purchases
\citep{koijen2021qe}, exchange rates \citep{koijen2020global},
benchmark-driven flows \citep{pavlova2023benchmarking}, and the
allocation of price-setting power across investor types
\citep{koijen2024which,haddad2025competitive}.

This note has two purposes. First, it lays out the workhorse logit
demand system and its identification (Part~I). Second, it confronts the
framework with two specific concerns that arise when one tries to use it
to model \emph{household} balance sheets across housing, deposits, and
equities as income grows (Part~II): (i)~whether the framework
accommodates \emph{non-homotheticity}---portfolio shares that vary
systematically with wealth or income---and (ii)~whether it accommodates
\emph{negative net portfolio shares}, of which the leading example is
the negative net housing position created by a mortgage. The household
case differs from the institutional case in a way that matters for
question~(ii): the negative leg (the mortgage) is bundled with the
positive leg (the house) by collateral and by a common price, so the
``shorts as supply'' fix that works for independent short books does
not transplant.

\bigskip
\hrule
\bigskip

\begin{center}\large\textsc{Part I: Institutional Demand-Based Asset Pricing}\end{center}

\section{The Logit Demand System for Portfolios}

Fix a date $t$ and let $i$ index investors and $n=1,\dots,N_t$ index
inside assets, with $n=0$ denoting an outside asset (cash, or, more
generally, the residual claim on non-asset wealth). Each investor has
assets under management $A_{i,t}$ and chooses portfolio weights
$\{w_{i,t}(n)\}_{n=0}^{N_t}$ summing to one. \citet{koijen2019demand}
specify the workhorse \emph{characteristics-based logit}:
\begin{align}
w_{i,t}(n) &= \frac{\delta_{i,t}(n)}{1 + \sum_{m=1}^{N_t}\delta_{i,t}(m)}, \qquad n=1,\dots,N_t, \label{eq:share}\\
w_{i,t}(0) &= \frac{1}{1 + \sum_{m=1}^{N_t}\delta_{i,t}(m)},\\
\delta_{i,t}(n) &= \exp\!\Bigl(\beta^0_{i,t}\,\me_{n,t} + \beta_{i,t}^\top x_{n,t} + \eps_{i,t}(n)\Bigr), \label{eq:delta}
\end{align}
where $\me_{n,t}$ is log market equity (or another scale-of-the-asset
variable), $x_{n,t}$ is a vector of observable asset characteristics
(profitability, dividends-to-book, betas, etc.), and $\eps_{i,t}(n)$ is
the latent demand for asset $n$ by investor $i$.

Three structural features of \eqref{eq:share} matter for what
follows. First, since $\delta_{i,t}(n) > 0$, every share $w_{i,t}(n) \in
(0,1)$: \emph{negative weights are ruled out by construction}. Second,
the system is Cobb-Douglas-like in the risky assets: relative weights
$w_{i,t}(n)/w_{i,t}(m)$ depend only on relative $\delta$'s, so they are
\emph{scale-free} in $A_{i,t}$. Third, market clearing is asset by
asset,
\begin{equation}
p_{n,t}\,s_{n,t} \;=\; \sum_i A_{i,t}\,w_{i,t}(n),
\label{eq:clearing}
\end{equation}
where $s_{n,t}$ is shares outstanding, so the equilibrium price
$p_{n,t}$ is implicitly defined and can be solved fixed-point style.
The map from demand parameters to equilibrium prices is what makes
counterfactuals (Was QE effective? Are markets inelastic? Who set the
price of Tesla?) tractable.

\section{Identification and Estimation}

The empirical challenge in \eqref{eq:delta} is that $\me_{n,t}$ and the
latent demand $\eps_{i,t}(n)$ are both functions of price and so
correlated: assets with high latent demand have high prices and so high
market equity. OLS estimates of $\beta^0_{i,t}$ are biased toward zero
(or even positive), implying spuriously inelastic demand.

\citet{koijen2019demand} resolve this with an instrument built from
investors' \emph{investment universes}. Let $\mathbf{1}_{i,t}(n) \in
\{0,1\}$ indicate whether asset $n$ is in investor $i$'s feasible set
at $t$ (estimated from the empirical extensive margin of holdings). The
instrument for $\me_{n,t}$ in investor $i$'s demand equation is the
counterfactual market equity that would obtain if every \emph{other}
investor $j \neq i$ held an equal-weighted portfolio over their own
universe:
\begin{equation}
\widehat{\me}_{n,t}^{\,(-i)} \;=\; \log\!\left( \sum_{j\neq i} A_{j,t}\,\frac{\mathbf{1}_{j,t}(n)}{\sum_m \mathbf{1}_{j,t}(m)} \right).
\end{equation}
The exclusion restriction is that other investors' AUM and universe
composition shift the supply of capital chasing asset $n$ without
shifting $i$'s latent demand for it. Conditional on asset
characteristics, this yields finite (and economically large) own-price
elasticities, typically near one---an order of magnitude below what
representative-agent calibrations would imply, and the empirical
foundation of the inelastic-markets literature
\citep{gabaix2021inelastic}.

The instrument has been the subject of two important refinements.
\citet{vanderbeck2022estimation} shows that when latent demand has a
common factor structure, the instrument is partly endogenous and
elasticities are biased; he proposes corrections that anchor the
estimated demand system to a small set of fundamentals.
\citet{haddad2025competitive} and \citet{koijen2024which} examine
heterogeneity in elasticities across investor types and find that a
small set of large, active investors do most of the price setting.

\section{Non-Homotheticity in the Institutional Setting}

The bare logit \eqref{eq:share} is homothetic in $A_{i,t}$ \emph{within
the risky set}: doubling AUM holding the $\delta$'s fixed scales every
position proportionally and leaves every share unchanged. To deliver
shares that move systematically with size, the framework must encode
size in either the coefficients or the characteristics. Four common
hooks:
\begin{enumerate}
\item Let $\beta^0_{i,t}$ depend on $\log A_{i,t}$, so the loading on
log market equity is size-elastic. Equivalently, allow the elasticity
of demand to vary with investor size.
\item Interact asset characteristics with investor characteristics:
$\beta_{i,t} = B \cdot z_{i,t}$ for a matrix $B$ and a vector
$z_{i,t}$ of investor attributes. \citet{koijen2024which} and
\citet{koijen2020global} pursue this aggressively, with separate
$\beta$'s estimated by investor type.
\item Allow the outside-asset margin $w_{i,t}(0)$ to depend on AUM
through a size-dependent constant in \eqref{eq:delta}.
\item Use random-coefficient (BLP-style) extensions in which the
distribution of $\beta_{i,t}$ over investors has wealth-varying
moments \citep{blp1995,nevo2001measuring}.
\end{enumerate}
All four are parametric: non-homotheticity is imposed through
functional form on the coefficients, not derived from preferences. The
distinction matters when extending to households (Section~\ref{sec:nh}
below), where preference-based theories of non-homotheticity already
exist \citep{wachter2010luxury}.

\section{Independent Short Selling}\label{sec:short-inst}

Because $\delta_{i,t}(n) > 0$, the logit cannot produce a negative
$w_{i,t}(n)$. In the original \citet{koijen2019demand} application this
is benign---the data are 13F long-only holdings, where shorts are
mostly absent. Four routes accommodate shorts when they matter:
\begin{enumerate}
\item \emph{Augmented choice set}: treat each asset as two products, a
long position and a short position, each with its own non-negative
logit share; net exposure is long minus short.
\item \emph{``Shorts as supply'' (twin logit)}: the investor brings a
long-side demand with budget $A^L_{i,t}$ and a short-side supply with
budget $A^S_{i,t}$, the latter bounded by the securities-lending
market \citep{davolio2002market,duffie2002securities,vayanos2008search,lewellen2011institutional}.
Market clearing becomes net, and signed exposure emerges in
equilibrium.
\item \emph{Mean-variance or linear demand}: naturally signed, at the
cost of losing the discrete-choice microfoundation and characteristics
parsimony.
\item \emph{Local linearisation around a no-short benchmark}:
elasticities are recovered analytically while estimation is restricted
to long holdings \citep{haddad2025competitive}.
\end{enumerate}
Crucially, each of these treats the short as \emph{independent} of any
particular long position: the investor freely allocates
shorting capacity across assets, subject to a separate budget. This is
the appropriate stylisation for an equity hedge fund, and it is not the
right stylisation for a mortgaged homeowner. Part~II addresses why.

\bigskip
\hrule
\bigskip

\begin{center}\large\textsc{Part II: Household Demand Systems with Leverage and Non-Homotheticity}\end{center}

\section{From Institutional to Household Demand Systems}

Transposing the logit framework to households changes several primitives.
The choice set narrows from thousands of securities to a handful of
asset \emph{classes}: housing, deposits, listed equities, bonds, and
private business equity. The outside-asset margin, a residual in the
institutional case, becomes a first-order modelling object---housing
dominates the median household balance sheet. Income, not just wealth,
is a structural driver of allocations, so the investor-characteristic
vector $z_{i,t}$ must include labour-income variables. Leverage is no
longer optional; many household balance sheets are built around a
mortgage. Finally, the relevant data are surveys (the U.S. Survey of
Consumer Finances, the euro-area Household Finance and Consumption
Survey) or administrative wealth registers (Sweden, Norway), not 13F.

The demand-system literature on households is younger than its
institutional counterpart but already extends to insurance choice
\citep{koijen2016health} and is being applied to balance-sheet
composition in ongoing work. Its appeal is the same as in the
institutional setting: characteristics-based parsimony, an explicit
demand function rather than a structural utility, and the prospect of
identifying elasticities from cross-sectional variation rather than
from contested time-series moments.

\section{Non-Homothetic Household Portfolio Shares}\label{sec:nh}

The household-finance literature documents that portfolio shares are
strongly non-homothetic. The risky share rises with wealth, both within
and across cohorts
\citep{calvet2007down,calvet2009fight,carroll2002portfolios,bach2020rich,fagereng2020heterogeneity};
gross housing's share of total assets declines with wealth above a
threshold, while deposits are hump-shaped or declining; mortgage debt
is hump-shaped over the life cycle. The \citet{gomes2021household}
survey collects this evidence.

Theoretically, two routes generate this empirically. The
\emph{structural} route derives non-homotheticity from preferences and
constraints: HARA utility, subsistence consumption, luxury goods, or
borrowing constraints that bind differently at different wealths
\citep{wachter2010luxury,cocco2005portfolio,yao2005optimal,chetty2017effect}.
The \emph{demand-system} route, in contrast, imposes
non-homotheticity through the parametric hooks listed in Section~4: let
$\beta_{i,t}$ in \eqref{eq:delta} depend on $(\log A_{i,t}, \log y_{i,t},
\mathrm{age}_{i,t})$ so that the loading on each asset characteristic
moves with wealth, income, and age. Concretely, for asset class
$n\in\{H,D,E\}$ (housing, deposits, equities),
\begin{equation}
\delta_{i,t}(n) \;=\; \exp\!\Bigl( \alpha_n + \beta^{w}_n\,\log A_{i,t} + \beta^{y}_n\,\log y_{i,t} + \beta^{a}_n\,\mathrm{age}_{i,t} + \gamma_n^\top x_{n,t} + \eps_{i,t}(n) \Bigr),
\label{eq:hh-delta}
\end{equation}
where the wealth-, income-, and age-dependent intercepts deliver
non-homothetic shares directly. Coefficients $\beta^w_n,\beta^y_n$ are
identified from the cross-section of households with different
$(A,y,\mathrm{age})$ holding the same characteristics $x_{n,t}$ fixed.

This route is flexible but \emph{descriptive}: \eqref{eq:hh-delta} can
reproduce any cross-sectional pattern of shares in $(A,y,\mathrm{age})$
but does not explain why. For applications where the question is
``what would happen to allocations if income $y$ moved by~$\Delta$,''
this is enough; for welfare or policy counterfactuals, one usually
wants a structural foundation alongside.

\section{Linked Positions: Housing and Mortgages}\label{sec:linked}

The user's motivating case---negative net housing shares created by
mortgages---fails the institutional fixes of Section~\ref{sec:short-inst}
for one reason. A mortgage is not an independent short. It is the
financing leg of a specific housing purchase, tied to the corresponding
gross housing position by a loan-to-value constraint
\begin{equation}
M_{i,t} \le \ell_{i,t}\,H_{i,t},
\label{eq:ltv}
\end{equation}
and by a common underlying price (the house price). Modelling a
``demand for mortgages'' as an independent supply, separately
characteristised and budgeted, severs that link. Four reformulations
preserve it.

\paragraph{Gross-portfolio reformulation with explicit LTV.}
The cleanest route, and the one we recommend for the user's
application. Run the logit \eqref{eq:share}--\eqref{eq:hh-delta} over
\emph{gross}, non-negative positions $\{H, D, E, B, M\}$, treating the
mortgage $M$ as a separate ``asset'' whose characteristic vector
includes a negative-yield component (the mortgage rate) and whose
choice is linked to $H$ via \eqref{eq:ltv}. Net portfolio shares are
then accounting identities:
\begin{equation}
\frac{H_{i,t} - M_{i,t}}{A_{i,t}}, \quad \frac{D_{i,t}}{A_{i,t}}, \quad \frac{E_{i,t}}{A_{i,t}}, \quad \dots,
\end{equation}
with $A_{i,t} = H_{i,t} + D_{i,t} + E_{i,t} + B_{i,t} - M_{i,t}$, and
they can take any sign. Negative net housing shares arise mechanically
for young leveraged households where $M_{i,t} > H_{i,t}$ (underwater)
or where $H_{i,t}$ enters $A_{i,t}$ in the denominator only through
$H - M$. Non-homothetic \emph{net} housing shares follow from
non-homothetic \emph{gross} $H$ and $M$ schedules in
\eqref{eq:hh-delta}, which scale differently in $(A,y,\mathrm{age})$.

The cost of this approach is that \eqref{eq:ltv} is an inequality
constraint, so closed-form market clearing breaks when it binds---one
solves the household problem subject to the constraint and aggregates
numerically. Identification of the elasticity of $M$ to the mortgage
rate exploits monetary-policy shocks and lender-market structure
\citep{greenwald2018mortgage,kaplan2020housing}.

\paragraph{Nested logit with continuous leverage.}
An alternative when leverage choice is continuous: an outer logit over
asset-class shares (Section~\ref{sec:nh}), then an inner Tobit /
fractional-response model for $\ell_{i,t}\in[0,\bar\ell]$ conditional
on $H_{i,t}>0$. This keeps the IO discrete-choice machinery for the
``what to hold'' decision and isolates the leverage choice for separate
identification (e.g., from regional variation in lender competition).

\paragraph{Bundle-as-product.}
If leverage is effectively discrete (a handful of mortgage products),
treat ``housing at LTV $\ell_k$'' as the $k$-th product entering the
logit alongside deposits and equities. Non-negativity is preserved at
the bundle level, and net housing share is mechanically
$(1-\ell_k)\,w_{H,\ell_k}$. The choice set grows, but the framework is
unchanged.

\paragraph{AIDS / QAIDS on net shares.}
The most direct way to allow signed and non-homothetic shares is to
abandon logit for an Almost Ideal Demand System
\citep{deaton1980almost} or its quadratic extension
\citep{banks1997quadratic}. These specify shares as linear (or
quadratic) functions of log-prices and log-total-expenditure, so they
are naturally signed and the ``Q'' term encodes non-homotheticity
directly. The trade-off is the loss of three things the logit gave us:
the McFadden microfoundation, characteristics-based parsimony (AIDS is
specified at the product level, not the characteristic level), and the
\citet{koijen2019demand} IV strategy. For exploratory empirical work on
household balance sheets, AIDS/QAIDS may in fact be the better starting
point.

\section{Mapping to Household Data}

The gross-portfolio reformulation aligns naturally with how household
data are organised. The SCF (US), the HFCS (euro area), and Swedish or
Norwegian administrative wealth registers all report assets gross of
liabilities (gross housing value, deposit balances, mutual-fund
holdings, business equity) and liabilities as separate line items
(primary-residence mortgage, other property loans, consumer credit,
student loans). The implementation maps directly:
\begin{itemize}
\item Gross $H,D,E,B,M$ from the wealth statement.
\item Investor characteristics $z_{i,t}$ from the demographic and
income blocks: age, education, family structure, labour income,
unemployment risk proxies.
\item Asset characteristics $x_{n,t}$: deposit and mortgage rates at
the local lender level, regional house-price-to-rent ratios, expected
equity-market characteristics (dividend yield, P/E) at the household's
investment horizon.
\end{itemize}
Identification is harder than in the institutional case. The
\citet{koijen2019demand} instrument relies on rich heterogeneity in
investment universes across investors; households' universes are
narrow and overlap heavily. Substitutes used in the household-finance
literature include: regional house-price shocks
\citep{kaplan2020housing}; monetary-policy-driven mortgage-rate
variation \citep{greenwald2018mortgage}; lender market structure;
cohort and age effects for life-cycle non-homotheticity; idiosyncratic
income shocks. The natural empirical strategy is to combine
characteristics-based demand estimation à la Koijen-Yogo with
quasi-experimental identification of a small number of key
elasticities, in particular the mortgage-rate elasticity of $M$
conditional on $H$.

Pitfalls worth flagging: top-coding of high wealth in surveys,
missingness in private business equity, mark-to-market versus
historical-cost housing values, mortgage-refinancing timing that
breaks contemporaneous identification of $M$ from current rates, and
the fact that LTV constraints \eqref{eq:ltv} bind selectively across
the wealth distribution.

\section{Concluding Remarks}

The logit demand framework of \citet{koijen2019demand} can be
transposed to household balance sheets but requires two extensions to
deliver what the user's application needs. Non-homotheticity (the
allocation pattern of housing, deposits, and equities as income grows)
is accommodated parametrically via wealth-, income-, and age-interacted
coefficients in \eqref{eq:hh-delta}---flexible but descriptive.
Negative \emph{net} portfolio shares (the negative net housing position
created by a mortgage) are accommodated naturally by a
gross-portfolio reformulation with an explicit LTV constraint, where
gross positions remain non-negative inside the logit and net shares are
accounting identities that can take any sign. The independent-short
fixes that work for institutional equity portfolios---augmented choice
set or ``shorts as supply''---do \emph{not} fit the household-leverage
case because the negative liability is bundled with a specific positive
asset.

For analysts willing to give up the discrete-choice microfoundation and
the Koijen-Yogo IV, an AIDS or QAIDS system on net shares is the
closer-fitting tool: it is naturally signed, naturally non-homothetic,
and has decades of empirical infrastructure behind it. The decision
between these two routes turns on whether the priority is characteristic-based
counterfactual analysis (favour the gross-portfolio logit) or flexible
description of the wealth-allocation profile (favour QAIDS).

\bibliographystyle{plainnat}
\bibliography{references}

\end{document}
